% Options for packages loaded elsewhere
\PassOptionsToPackage{unicode}{hyperref}
\PassOptionsToPackage{hyphens}{url}
\documentclass[
]{article}
\usepackage{xcolor}
\usepackage{amsmath,amssymb}
\setcounter{secnumdepth}{-\maxdimen} % remove section numbering
\usepackage{iftex}
\ifPDFTeX
  \usepackage[T1]{fontenc}
  \usepackage[utf8]{inputenc}
  \usepackage{textcomp} % provide euro and other symbols
\else % if luatex or xetex
  \usepackage{unicode-math} % this also loads fontspec
  \defaultfontfeatures{Scale=MatchLowercase}
  \defaultfontfeatures[\rmfamily]{Ligatures=TeX,Scale=1}
\fi
\usepackage{lmodern}
\ifPDFTeX\else
  % xetex/luatex font selection
\fi
% Use upquote if available, for straight quotes in verbatim environments
\IfFileExists{upquote.sty}{\usepackage{upquote}}{}
\IfFileExists{microtype.sty}{% use microtype if available
  \usepackage[]{microtype}
  \UseMicrotypeSet[protrusion]{basicmath} % disable protrusion for tt fonts
}{}
\makeatletter
\@ifundefined{KOMAClassName}{% if non-KOMA class
  \IfFileExists{parskip.sty}{%
    \usepackage{parskip}
  }{% else
    \setlength{\parindent}{0pt}
    \setlength{\parskip}{6pt plus 2pt minus 1pt}}
}{% if KOMA class
  \KOMAoptions{parskip=half}}
\makeatother
\usepackage{color}
\usepackage{fancyvrb}
\newcommand{\VerbBar}{|}
\newcommand{\VERB}{\Verb[commandchars=\\\{\}]}
\DefineVerbatimEnvironment{Highlighting}{Verbatim}{commandchars=\\\{\}}
% Add ',fontsize=\small' for more characters per line
\newenvironment{Shaded}{}{}
\newcommand{\AlertTok}[1]{\textcolor[rgb]{1.00,0.00,0.00}{\textbf{#1}}}
\newcommand{\AnnotationTok}[1]{\textcolor[rgb]{0.38,0.63,0.69}{\textbf{\textit{#1}}}}
\newcommand{\AttributeTok}[1]{\textcolor[rgb]{0.49,0.56,0.16}{#1}}
\newcommand{\BaseNTok}[1]{\textcolor[rgb]{0.25,0.63,0.44}{#1}}
\newcommand{\BuiltInTok}[1]{\textcolor[rgb]{0.00,0.50,0.00}{#1}}
\newcommand{\CharTok}[1]{\textcolor[rgb]{0.25,0.44,0.63}{#1}}
\newcommand{\CommentTok}[1]{\textcolor[rgb]{0.38,0.63,0.69}{\textit{#1}}}
\newcommand{\CommentVarTok}[1]{\textcolor[rgb]{0.38,0.63,0.69}{\textbf{\textit{#1}}}}
\newcommand{\ConstantTok}[1]{\textcolor[rgb]{0.53,0.00,0.00}{#1}}
\newcommand{\ControlFlowTok}[1]{\textcolor[rgb]{0.00,0.44,0.13}{\textbf{#1}}}
\newcommand{\DataTypeTok}[1]{\textcolor[rgb]{0.56,0.13,0.00}{#1}}
\newcommand{\DecValTok}[1]{\textcolor[rgb]{0.25,0.63,0.44}{#1}}
\newcommand{\DocumentationTok}[1]{\textcolor[rgb]{0.73,0.13,0.13}{\textit{#1}}}
\newcommand{\ErrorTok}[1]{\textcolor[rgb]{1.00,0.00,0.00}{\textbf{#1}}}
\newcommand{\ExtensionTok}[1]{#1}
\newcommand{\FloatTok}[1]{\textcolor[rgb]{0.25,0.63,0.44}{#1}}
\newcommand{\FunctionTok}[1]{\textcolor[rgb]{0.02,0.16,0.49}{#1}}
\newcommand{\ImportTok}[1]{\textcolor[rgb]{0.00,0.50,0.00}{\textbf{#1}}}
\newcommand{\InformationTok}[1]{\textcolor[rgb]{0.38,0.63,0.69}{\textbf{\textit{#1}}}}
\newcommand{\KeywordTok}[1]{\textcolor[rgb]{0.00,0.44,0.13}{\textbf{#1}}}
\newcommand{\NormalTok}[1]{#1}
\newcommand{\OperatorTok}[1]{\textcolor[rgb]{0.40,0.40,0.40}{#1}}
\newcommand{\OtherTok}[1]{\textcolor[rgb]{0.00,0.44,0.13}{#1}}
\newcommand{\PreprocessorTok}[1]{\textcolor[rgb]{0.74,0.48,0.00}{#1}}
\newcommand{\RegionMarkerTok}[1]{#1}
\newcommand{\SpecialCharTok}[1]{\textcolor[rgb]{0.25,0.44,0.63}{#1}}
\newcommand{\SpecialStringTok}[1]{\textcolor[rgb]{0.73,0.40,0.53}{#1}}
\newcommand{\StringTok}[1]{\textcolor[rgb]{0.25,0.44,0.63}{#1}}
\newcommand{\VariableTok}[1]{\textcolor[rgb]{0.10,0.09,0.49}{#1}}
\newcommand{\VerbatimStringTok}[1]{\textcolor[rgb]{0.25,0.44,0.63}{#1}}
\newcommand{\WarningTok}[1]{\textcolor[rgb]{0.38,0.63,0.69}{\textbf{\textit{#1}}}}
\usepackage{longtable,booktabs,array}
\usepackage{caption}
\captionsetup[table]{skip=6pt}
\newcounter{none} % for unnumbered tables
\usepackage{calc} % for calculating minipage widths
% Correct order of tables after \paragraph or \subparagraph
\usepackage{etoolbox}
\makeatletter
\patchcmd\longtable{\par}{\if@noskipsec\mbox{}\fi\par}{}{}
\makeatother
\setlength{\emergencystretch}{3em} % prevent overfull lines
\providecommand{\tightlist}{%
  \setlength{\itemsep}{0pt}\setlength{\parskip}{0pt}}
\usepackage{bookmark}
\IfFileExists{xurl.sty}{\usepackage{xurl}}{} % add URL line breaks if available
\urlstyle{same}
% fallback for those not using the hyperref driver hyperxmp:
\makeatletter
\@ifundefined{xmpquote}{\newcommand{\xmpquote}[1]{#1}}{}
\makeatother
\hypersetup{
  hidelinks,
  pdfcreator={LaTeX via pandoc}}

\author{}
\date{}

\begin{document}

{
\setcounter{tocdepth}{2}
\tableofcontents
}
\section{WebCoder --- Ghidul
administratorului}\label{webcoder--ghidul-administratorului}

\begin{quote}
Rezumat într-un paragraf: Administratorii rulează motorul de conținut
(cozi de curare, aprobări, rejudging, importuri) și motorul de siguranță
(moderare, raportări cu SLA, apeluri, scanări de integritate), plus
operațiunile platformei (status, backupuri, throttle-uri, curatorul-AI
opțional). Ghidul acoperă fiecare suprafață de admin și comenzile de
management din spatele ei.
\end{quote}

\textbf{Cui se adresează:} administratorii de platformă și VERIFIER-i.
Perechea în engleză:
\href{../guide-en/ADMIN-GUIDE.md}{\texttt{docs/guide-en/ADMIN-GUIDE.md}}.

\begin{center}\rule{0.5\linewidth}{0.5pt}\end{center}

\subsection{1. Dashboard-ul de admin}\label{1-dashboard-ul-de-admin}

\texttt{/admin/dashboard} --- contoare, plus tab-urile operaționale:
browser-ul de submissions (filtru pe verdict, buton \textbf{Rejudge}),
monitorul de submissions blocate (PENDING/COMPILING/RUNNING mai vechi de
N minute), managementul utilizatorilor (schimbările de rol --- singurul
loc unde se mișcă rolurile) și snapshot-ul de metrici
(\texttt{/api/v1/admin/metrics/}: adâncimea cozii, distribuția
verdicturilor, statisticile de 24 h ale duratelor de evaluare).

\subsection{2. Curarea conținutului}\label{2-curarea-conux21binutului}

\subsubsection{Coada de aprobare}\label{coada-de-aprobare}

\texttt{/admin/problem-queue} listează problemele
\texttt{status=PENDING} cu tot contextul de care are nevoie o evaluare:

\begin{itemize}
\tightlist
\item
  randarea enunțului + raportul linter-ului + scorul de lizibilitate +
  panoul de gap-uri de traducere;
\item
  \textbf{rândul Tests} --- perechile input/expected ale fiecărui test,
  fanionul demo, punctele, subtask-ul, bara per test, rata de prindere;
\item
  \textbf{puterea suitei de teste}
  (\texttt{/problems/\{id\}/test-strength/}) --- rata de prindere per
  test peste soluții cunoscut-ca-greșite, puterea suitei (excluzând
  exemplele), fanionul de necalibrare;
\item
  aprobare/respingere cu feedback (notifică autorul; respingerile
  păstrează problema editabilă).
\end{itemize}

\subsubsection{Rejudge și integritate}\label{rejudge-ux219i-integritate}

\begin{itemize}
\tightlist
\item
  \textbf{Rejudge} (admin pe toată platforma; VERIFIER pe problemele
  proprii): resetează la PENDING și re-dispatch-ează prin lane-ul real
  al evaluatorului; rezultatele per test de la prima execuție se
  persistă, deci swap-urile rapide de scor sunt posibile.
\item
  \textbf{Scanările de plagiat/boilerplate} --- badge-uri consultative
  pentru review uman, niciodată pedeapsă automată; vederea diff
  side-by-side; euristicile de multi-cont pe IP-uri alimentează aceeași
  coadă.
\item
  \textbf{Etichetarea soluțiilor-AI} --- review post-concurs al
  declarațiilor + anomaliilor de tipar de rezolvare (vezi coada valului
  de integritate).
\end{itemize}

\subsubsection{Importuri și operații de
catalog}\label{importuri-ux219i-operaux21bii-de-catalog}

Comenzi de management (rulate în containerul backend):

\begin{Shaded}
\begin{Highlighting}[]
\ExtensionTok{python}\NormalTok{ manage.py seed\_problems            }\CommentTok{\# pachetul de start + curriculum + taguri}
\ExtensionTok{python}\NormalTok{ manage.py boost\_hidden\_tests       }\CommentTok{\# regenerează suite ascunse din specificațiile comise}
\ExtensionTok{python}\NormalTok{ manage.py catalog\_polish           }\CommentTok{\# payload{-}ul de traduceri/igienă (idempotent)}
\ExtensionTok{python}\NormalTok{ manage.py catalog\_cleanup }\PreprocessorTok{[{-}{-}}\SpecialStringTok{apply}\PreprocessorTok{]}  \CommentTok{\# resturi de fixture e2e/ad{-}hoc}
\ExtensionTok{python}\NormalTok{ manage.py export\_catalog           }\CommentTok{\# snapshot determinist de fixture}
\ExtensionTok{python}\NormalTok{ manage.py import\_catalog           }\CommentTok{\# rebuild (sau {-}{-}from{-}release \textless{}tag\textgreater{})}
\ExtensionTok{python}\NormalTok{ manage.py calibrate\_difficulty     }\CommentTok{\# dificultate din ratele de AC (CNT{-}3)}
\end{Highlighting}
\end{Shaded}

Proveniența importurilor: ledger-ul CNT-2 \texttt{ImportSource} (sursă,
contact, termeni, URL de dovadă) trebuie să aibă un rând înainte ca un
import de corpus să fie legitim.

\subsection{3. Moderarea}\label{3-moderarea}

\begin{itemize}
\tightlist
\item
  \textbf{Coada de raportări}: fiecare buton de raportare (postări,
  mesaje, probleme, submissions) aterizează aici cu tracking de SLA;
  acțiuni: rezolvare/respingere cu notă.
\item
  \textbf{Apeluri}: deciziile de moderare sunt apelabile; audit trail-ul
  înregistrează operatorul, acțiunea și nota.
\item
  \textbf{Contact de abuz + pagina de bounty}: \texttt{/security}
  documentează fluxul de raportare; pagina de bug-bounty îl
  formalizează.
\item
  \textbf{Filtrul de conținut} (mesagerie): filtru RO+EN de tip
  marcare-nu-blocare --- mesajele marcate intră în coada de review, nu
  dispar silențios.
\item
  Ascunderea-la-prag + votul ponderat cu încredere rulează automat
  (negativele mici randează 0; auto-hide cu comutator de afișare).
\end{itemize}

\subsection{4. Curatorul-AI (opțional, IMPLICIT
OPRIT)}\label{4-curatorul-ai-opux21bional-implicit-oprit}

Evaluează propunerile de traducere ale comunității
(\texttt{/api/v1/translations/proposals/}). Configurare prin env
(\texttt{AI\_CURATOR\_URL} = orice endpoint compatibil OpenAI;
\texttt{AI\_CURATOR\_MODEL}; \texttt{AI\_CURATOR\_TIMEOUT\_S} ≤ 30;
\texttt{AI\_CURATOR\_AUTO\_APPROVE=false} implicit). Politica: trecerile
cu încredere ridicată intră în coada de confirmare umană; cea scăzută e
marcată; timeout-urile/erorile marchează \texttt{ai\_verdict=error} și
cad în coada umană --- fail-open, oamenii confirmă fiecare aprobare dacă
auto-aprobarea nu e explicit pornită. Oricare server compatibil OpenAI
funcționează (OpenAI, Azure, vLLM, Ollama\ldots). Vezi RUNBOOK
§ai-curator.

\subsection{5. Operațiunile
platformei}\label{5-operaux21biunile-platformei}

\begin{itemize}
\tightlist
\item
  \textbf{Pagina de status} (\texttt{/status}): adâncimea live a cozii,
  sănătatea evaluatorului; valul de ops a adăugat transparența cozii,
  scripturile de autoscale (\texttt{scripts/judge-autoscale.sh}) și
  harness-ul de load-test (\texttt{scripts/load\_smoke.py\ -\/-judge},
  \texttt{e2e/loadtest/replay.mjs}).
\item
  \textbf{Backupurile rulează implicit} (sidecar db-backup; RUNBOOK
  §Backups): declanșare manuală, listarea artefactelor, restore +
  verificare prin contoarele de rânduri din
  \texttt{/api/v1/users/admin/stats/}; script de exercițiu de restore
  trimestrial (\texttt{scripts/restore-drill.sh}).
\item
  \textbf{Throttle-urile} sunt reglabile prin env
  (\texttt{THROTTLE\_*}); pe redis.
\item
  \textbf{Deploy blue-green} + poarta de capacitate pentru concursuri
  (\texttt{manage.py\ contest\_replay\ -\/-contest\ \textless{}id\textgreater{}\ -\/-users\ 50})
  preced orice prim concurs cu rating.
\item
  \textbf{Email}: backend console implicit; setează \texttt{EMAIL\_HOST}
  + SPF/DKIM/DMARC înainte de trimiteri reale (RUNBOOK §email).
\end{itemize}

\subsection{6. Administrarea
utilizatorilor}\label{6-administrarea-utilizatorilor}

\begin{itemize}
\tightlist
\item
  Roluri: BASIC / CREATOR / TEACHER / VERIFIER / ADMIN --- se schimbă
  doar din vederea de admin a utilizatorilor (înregistrarea nu poate
  auto-atribui roluri).
\item
  2FA: enforce-ment progresiv pentru rolurile privilegiată
  (\texttt{TWO\_FACTOR\_ENFORCE} trece de la avertizare la obligație).
\item
  GDPR: export de date (\texttt{POST\ /users/me/export/}), soft-delete +
  purgare programată; profilurile minorilor ne-publice; datele de
  consimțământ parental.
\item
  Certificate: \texttt{manage.py\ issue\_certificates} (derivare
  idempotentă; \texttt{-\/-notify} opțional); ledger de revocare pentru
  cazurile de integritate.
\end{itemize}

\subsection{7. Rețete de incident}\label{7-reux21bete-de-incident}

{\def\LTcaptype{none} % do not increment counter
\begin{longtable}[]{@{}ll@{}}
\toprule\noalign{}
Simptom & Primele mișcări \\
\midrule\noalign{}
\endhead
\bottomrule\noalign{}
\endlastfoot
Submissions blocate în PENDING & monitorul de blocări → logurile
worker-ului (\texttt{docker\ compose\ logs\ celery-worker}) → rejudge
după fix \\
Verdict suspectat greșit & rejudge la submission; compară rândurile per
test; vezi \texttt{problems/\{id\}/test-strength/} \\
Poluare de catalog după rulări e2e &
\texttt{catalog\_cleanup\ -\/-apply\ -\/-allow-submissions} (gardată pe
conturi de probă) \\
Beat nu programează & \texttt{docker\ compose\ logs\ celery-beat} ---
flagul e \texttt{-\/-schedule} (celery 5.x) \\
Backend blocat & \texttt{/healthz} vs \texttt{/readyz} disting
liveness-ul de dependențe; restart-ul unless-stopped se auto-vindecă \\
\end{longtable}
}

\end{document}
